\section{Evaluation}
\label{sec:eval}

Five questions drive the evaluation:

\begin{description}
\itemsep0pt
\item[Q1] What does the gate cost per message? (\Cref{sec:eval:rtt})
\item[Q2] Is the slot-indexed validate really $O(1)$? (\Cref{sec:eval:o1})
\item[Q3] Does the gated path hold up under MPSC contention?
          (\Cref{sec:eval:mpsc})
\item[Q4] How quickly does revocation take effect?
          (\Cref{sec:eval:revoke})
\item[Q5] What does the gate touch in the cache hierarchy?
          (\Cref{sec:eval:perfctrs})
\end{description}

\subsection{Hardware and methodology}
\label{sec:eval:hw}

All measurements were taken on an Intel Core i7-11370H (3.30\,GHz,
Tiger Lake, 4 cores / 8 threads) running Fedora 43 with kernel 7.0.10.
The \texttt{performance} cpufreq governor is active, and the kernel
command line carries
\texttt{nosmt isolcpus=2,3 nohz\_full=2,3 rcu\_nocbs=2,3}---SMT off,
two physical cores reserved---with producer and consumer threads
pinned to the isolated pair. We disable SMT by kernel parameter
rather than in firmware so that the configuration is fully captured
by \texttt{/proc/cmdline}, and therefore reproducible by anyone
running the artefact. The toolchain is
C++23 on gcc 15.2.1 at \texttt{-O3 -march=native}. Time comes from
\texttt{rdtscp + lfence}, calibrated against \texttt{steady\_clock}
over a 50\,ms window.

Both of those instructions serialise, so the bracket is not free, and we
measure its cost rather than assume it away: an empty
\texttt{now()}-to-\texttt{now()} pair takes \timerFloor\,ns on this
host. Every absolute latency below carries that as an additive constant,
and we publish it so a reader can subtract it. It cancels in the gate
cost, which is a difference between two equally bracketed measurements,
but it does not cancel in the absolutes---and because it is a constant,
it inflates the fastest transports proportionally most.

We note the platform explicitly because it shapes what these numbers
can support. This is a mobile-class part, not a server: absolute
latencies would differ on server silicon, and sustained load invites
thermal effects that a rack part would not see. The claims we defend
are therefore \emph{relative}---gate-on versus gate-off in the same
binary on the same machine---and \emph{shape}, such as flatness of
\texttt{validate} across pool occupancy. Both are far less sensitive
to the platform than any absolute floor, and the artefact reruns
everywhere.

\textbf{Warm-up} is 1,000,000 untimed round trips per measurement.
The count is large deliberately. Warm-up is a shared iteration count
rather than a duration, so it buys the \emph{fastest} transport the
\emph{least} settling time; at a smaller count our lowest-latency
paths measured before the frequency ramp had settled, which showed up
as a systematic drift across repetitions on exactly those rows.
\textbf{Iteration count} is 200,000 per data point.
\textbf{Percentiles} come from sort-and-index over the full sample,
with no HDR-bucket approximation; the artefact ships every
per-repetition CSV alongside the merged medians.
\textbf{Repetitions.} Every configuration runs fifteen times end to
end, each in a fresh process, and the repetitions are
\emph{interleaved} rather than blocked: repetition $k$ of every
transport runs before repetition $k{+}1$ of any of them. That ordering
is not cosmetic. The gate cost is a difference between two transports,
so measuring all of one and then all of the other lets any drift across
the sweep land directly in that difference; interleaving makes drift
common-mode, so it cancels. We report the median run, and the
cross-repetition coefficient of variation per column rather than as one
aggregate, for the reason given in \Cref{sec:eval:rtt}. The sweep driver records the host environment---kernel,
governor, isolation flags, SMT state, compiler---next to every CSV.

\subsection{Q1: Per-message gate cost}
\label{sec:eval:rtt}

\Cref{fig:figure-1} plots p99 round-trip latency against payload size
for every transport in the comparison set; \Cref{tab:rtt} gives the
numbers at the canonical 256\,B payload.

\begin{figure}[t]
\centering
\includegraphics[width=\columnwidth]{figures/paper1_figure_1}
\caption{Round-trip p99 latency vs payload size, log-log. \sysname
gate ON (solid orange) and gate OFF (solid green) track each other
at a near-constant offset---the gate cost---across the payload
range, and both sit below every other transport measured.}
\label{fig:figure-1}
\end{figure}

\begin{table}[t]
\centering\footnotesize
\setlength{\tabcolsep}{4pt}
% Placeholder for Table 1 — generated by scripts/plot_paper1_figure.py
% with --tex flag during `./build.sh --refresh`.
\textit{[Table 1 not yet rendered. Run \texttt{./build.sh --refresh}
on a Linux x86\_64 host after \texttt{./scripts/run\_paper1\_sweep.sh}.]}

\caption{Round-trip p50 / p99 / p99.99 latency at 256\,B payload,
sorted by p99. Auto-generated by \texttt{plot\_paper1\_figure.py}.
Aeron is the one row measured without core isolation, for the reason
given in \Cref{sec:eval:threats}; its p99.99 is host scheduling on a
four-core mobile part rather than a property of Aeron, and we draw no
conclusion from it.}
\label{tab:rtt}
\end{table}

\textbf{Gate cost.} Each round trip validates twice---\IPCSEND at the
producer, \IPCRECV at the consumer---so the per-message cost is half
the gated--raw delta:
\begin{equation*}
\text{gate cost}_{p99} \;=\;
\bigl(\text{RTT}_{\text{gated}} - \text{RTT}_{\text{raw}}\bigr) / 2
\;=\; \gateOverheadTail\ \text{ns}
\end{equation*}
(RTT values as in \Cref{tab:rtt}: gated \gatedAstraRTT\,ns, raw
\rawAstraRTT\,ns.)
\textbf{On stability.} A single coefficient-of-variation count across
all merged cells would be misleading, so we report it per column.
Every \texttt{max} cell exceeds a 5\% threshold, which is arithmetic
rather than a finding: a maximum is one sample per repetition and has
no variance to average away. The p99.99 column behaves similarly for
the same reason. In the two columns this paper actually quotes, p50 and
p99, the overwhelming majority of cells fall inside 5\%. The tail
columns have a second problem besides being near-single samples: the
timing instrument's own p99.99 is \timerFloorTail\,ns, which is larger
than the entire median round trip of our fastest transports. Their
p99.99 figures therefore describe the instrument at least as much as the
transport, and we draw no conclusions from them. We give the
breakdown rather than the total because the total invites a conclusion
the data does not support, and the artefact ships the full per-cell
report either way.

The gate cost inherits the noise of both measurements it differences,
so its per-repetition values span a wider relative range than either
endpoint does. Bootstrapping the median across the repetitions puts its
95\% confidence interval at \gateCostCILo--\gateCostCIHi\,ns. We give
the interval rather than the point estimate alone because it straddles a
rounding boundary, and a bare integer would imply a precision the data
does not support. What the data does support is the claim that matters:
every value in that interval sits two orders of magnitude inside the
50\,ns design target.

\textbf{Relative cost.} The gate adds \gateOverheadPercent\% to the
round trip as measured. We quote that fraction ourselves because it is
the number a sceptical reader will compute, and because the honest
reading of it is favourable: a large percentage of \rawAstraRTT\,ns is
still a small number of nanoseconds, and the mediated path at
\gatedAstraRTT\,ns remains faster than every transport in
\Cref{tab:rtt}.

The fraction is also conservative. Both round trips carry the
\timerFloor\,ns instrument constant, and subtracting it from each
shrinks the denominator while leaving the difference untouched, so the
gate's true share of an uninstrumented round trip is somewhat higher
than what we print. We report the measured ratio anyway, because it is
what a reader can derive from \Cref{tab:rtt} without trusting our
arithmetic, and we would rather a correction move against us than for
us.

\textbf{On the Aeron comparison.} Aeron is in the table as a sanity
check on our substrate, not as a scalp. Its purpose is to establish
that the fastpath we are adding a gate to is a credible one rather
than a strawman---if our ungated numbers were mediocre, a cheap gate
would prove nothing. Aeron's published $\sim$250\,ns figure is taken at
a 100\,B payload, and our measured p50 brackets it: we land below it at
64\,B and above it at 256\,B (\Cref{fig:figure-1}), which is where a
100\,B point belongs. Matching a published number we did not produce,
on hardware that is not theirs, is what tells us the harness is
faithful. Its p99 on this platform is a different story---dominated by
scheduling jitter that a tuned server would not exhibit---so we do not
lean on it. That \sysname measures faster
still is unsurprising and not a claim we press: Aeron routes every
message through a separate media-driver process and carries flow
control, loss recovery, and a network transport that we simply do not
implement. The comparison establishes that our fastpath is fast. It
does not establish that \sysname is a better system than Aeron, and
we make no such argument (\Cref{sec:disc:aeron}).

\subsection{Q2: $O(1)$ validate across pool occupancy}
\label{sec:eval:o1}

\Cref{fig:figure-2} plots per-call validate latency (p50/p99/p99.99)
as pool occupancy grows from 16 to 4090 active tokens. The line is
flat to within measurement noise: the slot-indexed lookup simply does
not see pool size.

\begin{figure}[t]
\centering
\includegraphics[width=\columnwidth]{figures/paper1_figure_2}
\caption{validate() per-call latency vs token-pool occupancy. The
flat line is the $O(1)$ claim.}
\label{fig:figure-2}
\end{figure}

For contrast, the $O(n)$ pool scan this design replaced performs one
constant-time compare per live token---up to 4096 at the default pool
capacity---where the slot-indexed path performs exactly one after a
shift and a bounds check (\Cref{sec:design:fastpath}). We give that
contrast asymptotically rather than measuring it: the scan is not code
the artefact ships, and a structural difference a reader can verify
from \Cref{sec:design:fastpath} is worth more than a latency no
reviewer could reproduce. What we do measure is its consequence
(\Cref{fig:figure-2})---the same \validateTail\,ns whether 16 tokens
are live or 4090.

\subsection{Q3: MPSC contention scaling}
\label{sec:eval:mpsc}

This is the one question this host cannot answer, and we would rather
say so than fill the gap. The MPSC harness runs one thread per producer
plus a consumer and does not pin, so those threads take whatever the
general scheduler still owns---which \texttt{isolcpus} deliberately
reduces to two cores here. Every producer count above one therefore
oversubscribes, and the rows past that point describe thread contention
rather than the ring. The signature is unmistakable: aggregate
throughput falls by nearly three orders of magnitude between two and
four producers, and at eight the gated arm appears to overtake the
ungated one, which cannot happen for a gate that only adds work.

The single unoversubscribed point is worth reporting. With one producer
and one consumer, enabling the gate costs \mpscGateCostPercent\% of
aggregate throughput---consistent with the per-message cost in
\Cref{sec:eval:rtt}, and the only MPSC figure we defend. The scaling
sweep needs a host with cores to spare; the harness ships with the
artefact and runs there unmodified, and \Cref{sec:eval:threats} records
the limitation.

\begin{figure}[t]
\centering
\includegraphics[width=\columnwidth]{figures/paper1_figure_3_mpsc}
\caption{MPSC aggregate throughput vs producer count. Only the
single-producer point is free of oversubscription on this host: beyond
it the harness needs more runnable threads than the scheduler has cores
left after \texttt{isolcpus}, and the cliff is that contention rather
than any property of the ring. Included because the artefact reproduces
it, not because we draw a scaling conclusion from it.}
\label{fig:mpsc}
\end{figure}

\subsection{Q4: Revocation latency}
\label{sec:eval:revoke}

\Cref{fig:revoke} shows the empirical CDF of the window from the
\texttt{revoke()} call to the first \texttt{PERMISSION\_DENIED} seen
by a tight-loop producer, across pool occupancies
$\{16, 256, 2048, 4000\}$. The p99 is \revokeTail\,$\mu$s at every
pool size, and the tail is dominated by the producer's own loop
interval rather than by the cascade.

\begin{figure}[t]
\centering
\includegraphics[width=\columnwidth]{figures/paper1_figure_4_revoke}
\caption{Revocation-latency CDF. The producer keeps writing until it
sees its first \texttt{PERMISSION\_DENIED}; the window is measured
from the moment the reaper issued \texttt{revoke()}.}
\label{fig:revoke}
\end{figure}

\subsection{Q5: Microarchitectural footprint}
\label{sec:eval:perfctrs}

\Cref{tab:perf} reports cycles, L1d misses, LLC misses, and branch
mispredictions per \texttt{validate()}, collected with grouped
\texttt{perf\_event\_open} counters across pool occupancies. The
counts stay flat as the pool grows---consistent with $O(1)$. The L1d
budget is dominated by the slot load, and LLC misses stay in single
digits per call, so the hot slot effectively lives in L2 under
realistic workloads.

\begin{table}[t]
\centering\small
% Placeholder for the perf-counter table — generated by
% scripts/gen_table_perf.py during `./build.sh --refresh`.
\textit{[Performance-counter table not yet rendered. Run
\texttt{./build.sh --refresh} on a Linux x86\_64 host with
\texttt{perf\_event\_paranoid} $\le$ 1.]}

\caption{Microarchitectural cost of validate() per call.
Auto-generated from \texttt{paper1\_perfcounters.csv}.}
\label{tab:perf}
\end{table}

\subsection{Functional correctness}
\label{sec:eval:correctness}

\Cref{tab:tests} maps the correctness invariants we claim to their
runnable test artefacts. All of them pass on the
artefact-evaluation machine.

\begin{table}[t]
\centering\scriptsize
\setlength{\tabcolsep}{3pt}
\begin{tabular}{ll}
\toprule
Invariant                          & Test \\
\midrule
T-Forge --- forged slot rejected     & \texttt{test\_capability} (T5) \\
T-OOB --- out-of-range slot          & \texttt{test\_capability} (T6) \\
T-Mono --- superset derive fails      & \texttt{test\_capability} (T2) \\
T-Cascade --- descendant revoke       & \texttt{test\_capability} (T4) \\
T-RingFP --- denied = zero footprint  & \texttt{test\_ipc\_sprint5\_cap\_gate} (T7) \\
T-Conc --- no stale validate         & \texttt{test\_capability\_concurrency} \\
T-Surv --- revoke under load         & \texttt{test\_capability\_survivability} \\
T-Prop --- oracle agreement          & \texttt{test\_capability\_property} \\
T-Mut --- gate under MPSC + revoke   & \texttt{test\_ipc\_sprint5\_gate\_concurrency} \\
T-CBMC --- monotonicity (37 checks)  & \texttt{formal/cbmc/} \\
\bottomrule
\end{tabular}
\caption{Correctness invariants and their test artefacts. The
property test runs 50,000 operations $\times$ 5 seeds against a
hand-written C++ reference oracle.}
\label{tab:tests}
\end{table}

\subsection{Threats to validity}
\label{sec:eval:threats}

\para{Single microarchitecture} Every number comes from one
mobile-class Tiger Lake part (Core i7-11370H, \Cref{sec:eval:hw}).
Absolute latencies will move on other parts---server Intel,
AMD, aarch64---and sustained load on a laptop part invites thermal
effects a rack machine would not show. The claims we defend are
\emph{relative} (gate on vs.\ gate off, same binary, same machine)
and \emph{shape} (flatness across pool occupancy), both far less
sensitive to the part than any absolute floor. The artefact reruns
everywhere.

\para{Aeron comparison} Aeron's 250\,ns figure is its published
number on its authors' hardware. Our primary comparison is our own
Aeron measurement, taken in the same TSC domain as every other
transport (\texttt{bench\_baseline\_aeron}); the published floor is
context. If our local Aeron build underperforms its published number,
we report both rather than choosing the flattering one.

\para{Additive gate-cost model} Deriving gate cost as
$(\text{gated}-\text{raw})/2$ assumes the gate composes additively
with ring costs. The perf-counter profile
(\Cref{sec:eval:perfctrs}) supports that reading---the gate adds a
bounded number of loads and no branches on the accept path---but any
microarchitectural interaction, say the slot load evicting a ring
line, is folded into the delta rather than isolated.

\para{Timer error} Calibrating RDTSC to nanoseconds against
\texttt{steady\_clock} over a 50\,ms window bounds calibration error
well under 1\%, and the \texttt{rdtscp+lfence} serialisation cost is
identical in every transport's measurement loop, so it cancels in
comparisons. It does not cancel in absolute figures, which is why we
measure it (\timerFloor\,ns, \Cref{sec:eval:hw}) instead of only
asserting that it cancels. A per-call figure of the same order as the
instrument should be read as an upper bound, not a floor.

\para{Aeron runs without core isolation} Every other harness pins its
threads to the isolated cores; Aeron cannot. Its client adds a
conductor thread to our measuring and echo threads, so three threads
need CPU while two of them busy-spin, and two isolated cores with
\texttt{nohz\_full}---whose whole purpose is to stop the timer tick
there---leave no preemption to let the third run. The harness wedges.
Aeron therefore runs across all CPUs, which gives it a noisier tail
than the isolated measurements beside it. Its cross-repetition
behaviour shows the same thing: over fifteen repetitions its p99 at
256\,B moved by roughly a third between the best and worst run, a
spread no isolated harness came close to. We flag both rather than
bury them, because together they are the reason to read Aeron's tail
\emph{and} its run-to-run variation on this platform as properties of
the configuration we were forced into, and to lean only on the
substrate argument of \Cref{sec:eval:rtt}.

\para{MPSC scaling is unmeasured} The MPSC harness needs one runnable
thread per producer plus a consumer, and it does not pin, so it competes
for the cores \texttt{isolcpus} leaves to the general scheduler---two,
on this machine. Only the single-producer configuration fits. We report
that point and decline to draw a scaling curve from the rest, because
those rows measure our own thread oversubscription. This is a property
of the host, not of the design, and it is the one claim in the paper
that a machine with spare cores would straightforwardly settle.

\para{io\_uring configuration} Our io\_uring harness carries the
payload over a pipe and issues one \texttt{IORING\_OP\_WRITE} or
\texttt{IORING\_OP\_READ} per round trip, waiting on its completion
before submitting the next. That is io\_uring used as a syscall
replacement, not as the batching interface it exists to be, so this
column is close to the pipe column plus submission overhead---which is
what we measure, and it is not io\_uring at its best. A request/response
round trip has no independent work to batch, which is why the harness
is built this way; a reader should nonetheless not read this row as
io\_uring's floor.

We also report submit-per-operation rather than submission-queue
polling (\texttt{IORING\_SETUP\_SQPOLL}), and that reason is a property
of our platform rather than a preference. SQPOLL trades the submit syscall
for a kernel poller thread dedicated to each ring, and our
request/response harness opens two; on a four-core machine with two
cores removed by \texttt{isolcpus}, those pollers compete with the
benchmark itself. Measured with SQPOLL enabled, io\_uring was
\emph{slower} than without it. A reader should therefore not take our
io\_uring column as io\_uring's floor---on a host with cores to spare
SQPOLL is the right configuration, and the artefact enables it via
\texttt{ASTRA\_IOURING\_SQPOLL=1}. The comparison we actually depend
on is gate-on versus gate-off, which this does not touch.
